% Generated draft: requires mathematical and source review.
\documentclass[11pt]{article}
\usepackage[margin=25mm]{geometry}
\usepackage{fontspec}
\setmainfont{Latin Modern Roman}
\usepackage{amsmath,amssymb,mathtools,unicode-math}
\setmathfont{Latin Modern Math}
\usepackage{enumitem,needspace,xurl,hyperref}
\setlength{\parindent}{0pt}
\setlength{\parskip}{5pt}
\setlength{\emergencystretch}{3em}
\newcommand{\N}{\mathbb{N}}
\newcommand{\Z}{\mathbb{Z}}
\newcommand{\Q}{\mathbb{Q}}
\newcommand{\R}{\mathbb{R}}
\newcommand{\C}{\mathbb{C}}
\newcommand{\F}{\mathbb{F}}
\newcommand{\A}{\mathbb{A}}
\newcommand{\PP}{\mathbb{P}}
\newcommand{\E}{\mathbb{E}}
\newcommand{\eps}{\varepsilon}
\newcommand{\dd}{\,\mathrm{d}}
\begin{document}
\section*{UnsolvedMath ID: 30005604}
\textbf{Source problem code:} catalogue entry
\subsection{Definitions and mathematical statement}
The supplied record concerns Fourier-mode solutions of inhomogeneous Laplace equations expressed through modified Bessel functions. For \(\eta\in\C\), the symbols \(I_\eta\) and \(K_\eta\) denote the modified Bessel functions of the first and second kind, respectively. Their arguments and any required branch conventions are not supplied.

A Laurent polynomial in \(y\), over a coefficient field \(F\), is a finite sum
\[
p(y)=\sum_{k=m}^{M}c_k y^k,\qquad m,M\in\Z,\quad m\le M,\quad c_k\in F.
\]
This definition does not specify the coefficient field for the problem. Negative powers require \(y\ne0\) when the polynomial is evaluated. The original fragment describes quantities denoted by \(q_{i,j}\) and \(\mu_{i,j}\) as Laurent polynomials in \(y\), apparently depending on labels \(n_1,n_2,\lambda,\alpha,\beta\). The index ranges, parameter domains, and precise placement of these labels are not recoverable from the supplied fragment.

The requested statement, with its unresolved qualifications explicit, is: for every admissible tuple \((r,a,b)\), establish the existence of Fourier-mode solutions expressible through modified Bessel functions, obtain closed formulas for their Laurent-polynomial coefficients, and establish a connection between the relevant solution space and differential Galois theory.

This is an incomplete specification of a mathematical problem. Neither the set of admissible tuples nor the inhomogeneous Laplace equation is supplied. The underlying space, operator, forcing term, Fourier decomposition, Bessel-function ansatz, and boundary or growth conditions are absent. The relation between \((r,a,b)\) and \((n_1,n_2,\lambda,\alpha,\beta)\) is also unspecified. Consequently, the universal quantifier cannot be assigned a precise domain, and the existence and coefficient-formula requests cannot be made into a definite theorem from this record alone.

The original fragment explicitly proposes establishing a connection to differential Galois theory as an open problem, but does not specify the differential equation, base differential field, derivation, associated homogeneous equation, Galois group, or property to be related. Thus this part remains a research objective rather than a uniquely determined mathematical assertion. The fragment does not establish whether the existence and coefficient-formula requests in the reconstructed statement were themselves posed as open problems.
\subsection{Short English statement}
For every admissible parameter tuple, the reconstructed statement asks for Fourier-mode solutions involving modified Bessel functions, explicit Laurent-polynomial coefficients, and a connection to differential Galois theory. The admissibility conditions, equation, solution conditions, and intended Galois-theoretic conclusion are missing. The original fragment explicitly identifies only the proposed Galois-theoretic connection as an open problem.
\subsection{Sources}
\begin{itemize}
\item Supplied local UnsolvedMath registry record, \texttt{data/unsolvedmath/problems.json}, permanent ID \(30005604\), title: Bessel-Function Solutions of Inhomogeneous Laplace Equations. Exact source problem code: \texttt{catalogue entry}. Collection: Oberwolfach Reports Open Problems. No individual proposer is identified; the supplied proposed year is 2023.
\item Supplied report citation: Geometric Spectral Theory (2024), Oberwolfach Reports, DOI 10.4171/owr/2023/36, pp.\ 2017--2106. Supplied source link: \url{https://doi.org/10.4171/owr/2023/36}. The record describes Oberwolfach Reports as documenting mathematical workshops and their open-problem discussions; no narrower passage location is supplied.
\item Supplied literature source: Oberwolfach Report 36, \url{https://doi.org/10.4171/owr/2023/36}. Literature checked at: \texttt{2026-08-22}. This is supplied dated curation, not an independently verified assessment. It labels the problem open and reports that the proposed connection between the constructed Bessel-type equations and differential Galois theory is explicitly open, with no complete treatment located. The supplied background requests independent verification of time-sensitive status and cited claims.
\item The registry labels the reconstructed statement \texttt{corrected\_verified} and states that it was checked against the official report. These are supplied provenance claims, not verification performed for this draft. The supplied original text is a fragment and includes a footnote stating that a Python program used integer arithmetic only; the associated computation is not supplied.
\end{itemize}
\end{document}
